\chapter{Sylow Theorems and Finite Groups}

\begin{center}\itshape
Is there a simple group of order $30$?\\
Lagrange says what subgroup orders are allowed; Sylow says which ones must occur, and how many times.\\
Counting them answers the question in three lines.
\end{center}

Can a group of order $30$ be simple? Lagrange's theorem says that a subgroup has an order dividing $30$. It does not tell us which subgroups must exist.

The first improvement is Cauchy's theorem: because $2$, $3$ and $5$ divide $30$, the group contains an element of each of those orders. Sylow's theorems go further. They promise a subgroup whose order is the largest power of any chosen prime dividing the group, and they control how many such subgroups there can be.

For order $30$, the number $n_5$ of subgroups of order $5$ must divide $6$ and be congruent to $1$ modulo $5$. Thus $n_5$ is $1$ or $6$. Similarly $n_3$ is $1$ or $10$. If either count is $1$, that subgroup is normal and the group is not simple. If both counts take their larger values, the six subgroups of order $5$ contribute $24$ different non-identity elements, and the ten subgroups of order $3$ contribute another $20$. That is too many for a group of order $30$. No simple group of that order exists.

The counting argument is short once the Sylow theorems are available. The proofs of those theorems require a careful use of group actions, which is why we studied orbits first.

At the end we ask a second question: is there a finite chain of subgroups, each normal in the next, with abelian quotients? This gives the definition of a \emph{solvable group}. Galois theory will connect that property with solving polynomial equations by radicals.

\section{The three theorems}

Throughout the Sylow statements, $G$ is finite, $p$ is prime, and
$|G|=p^a m$ with integers $a\ge0$, $m\ge1$ and $p\nmid m$.
Thus $p^a$ is the largest power of $p$ dividing $|G|$; it is $1$
when $p$ does not divide $|G|$.

\begin{definition}
A \emph{$p$-subgroup} of $G$ is a subgroup whose order is a power of $p$. A \emph{Sylow $p$-subgroup} is a subgroup of order $p^a$. Their set is $\operatorname{Syl}_p(G)$ and their number $n_p$. The \emph{normaliser} of a subgroup $H$ is $N_G(H)=\{g\in G:gHg^{-1}=H\}$.
\end{definition}

$N_G(H)$ is the stabiliser of $H$ for the action of $G$ on its subgroups by conjugation, so it is a subgroup, and $H\trianglelefteq N_G(H)$; it is the largest subgroup of $G$ in which $H$ is normal. By the orbit--stabiliser theorem, $H$ has $[G:N_G(H)]$ conjugates.

\begin{lemma}\label{lem:binom}
Let $p$ be prime, $a\ge0$, and $m\ge1$. Then
\[
\binom{p^am}{p^a}\equiv m\pmod p.
\]
In particular, if $p\nmid m$ then this binomial coefficient is not divisible by $p$.
\end{lemma}
\begin{proof}
If $a=0$, then $\binom{m}{1}=m$, so both sides of the congruence are $m$. The rest of the argument is the case $a\ge1$.

We compare coefficients modulo $p$, that is, in polynomials with
coefficients in $\mathbb F_p=\mathbb Z/p\mathbb Z$. Here $p\cdot1=0$,
and no smaller positive multiple of $1$ is zero; this property is called
\emph{characteristic $p$}. Reducing integer coefficients modulo $p$
respects addition and multiplication, so we may reduce a polynomial
identity coefficient by coefficient.

First, for a prime $p$,
\[
(1+X)^p\equiv1+X^p\pmod p.
\]
Indeed the intermediate coefficient $\binom pj$
is divisible by $p$ whenever $0<j<p$.
Raise this identity to the $p$-th power
repeatedly. At each stage we use
$(U+V)^p=U^p+V^p$ in characteristic $p$,
and obtain
\[
(1+X)^{p^a}=1+X^{p^a}
\quad\text{in }\mathbb F_p[X].
\]

Now raise both sides to the power $m$:
\[
(1+X)^{p^am}
=(1+X^{p^a})^m
=\sum_{j=0}^m\binom mjX^{jp^a}
\quad\text{in }\mathbb F_p[X].
\]
Compare the coefficients of $X^{p^a}$.
On the left it is $\binom{p^am}{p^a}$;
on the right only $j=1$ contributes, giving
$\binom m1=m$. Therefore
\[
\boxed{\binom{p^am}{p^a}\equiv m\pmod p}.
\]

This congruence holds for every positive integer $m$.
We can conclude that the coefficient is \emph{non-zero}
modulo $p$ only when $p\nmid m$, which is
the case needed in the Sylow argument.
For a small check, with $p=2$, $a=1$, $m=3$,
we have $\binom62=15\equiv1\equiv3\pmod2$.
\end{proof}

\begin{theorem}[Sylow I]\label{thm:sylow1}
If $G$ is finite and $|G|=p^a m$ with $p\nmid m$, then $G$ has a subgroup of order $p^a$.
\end{theorem}
\begin{proof}[Proof (after Wielandt, 1959)]
Write $|G|=p^a m$, with $p\nmid m$. Let $\Omega$ be the collection of all subsets of $G$ containing exactly $p^a$ elements. The group acts on these subsets by left multiplication:
\[
g\cdot A=\{ga:a\in A\}.
\]
This really is an action: $eA=A$ and $(gh)A=g(hA)$.

The number of such subsets is $\binom{p^am}{p^a}$, which is not divisible by $p$ by Lemma~\ref{lem:binom}. Split $\Omega$ into orbits. If every orbit had size divisible by $p$, then their sum $|\Omega|$ would be divisible by $p$ too. Hence some orbit $O$ has size prime to $p$.

Choose a subset $A\in O$ and let $H=G_A$ be its stabiliser. Orbit--stabiliser now gives
\[
|G|=|O|\,|H|.
\]
Since $p$ does not divide $|O|$, the factor $p^a$ in $|G|$ must divide $|H|$. This is our lower bound.

For the upper bound, fix any $x\in A$. Every $h\in H$ sends $A$ to itself, so $hx\in A$. The map
\[
H\longrightarrow A,\qquad h\longmapsto hx
\]
is injective: if $h_1x=h_2x$, multiplying on the right by $x^{-1}$ gives $h_1=h_2$. Thus $|H|\le|A|=p^a$.

We have found a subgroup $H$ whose order is both a multiple of $p^a$ and at most $p^a$. Therefore $|H|=p^a$.
\end{proof}

\paragraph{Why this works.}
An action is \emph{free} when no element other than the identity fixes a point. There are two actions in the proof. The action of $G$ on the collection $\Omega$ of subsets need not be free; its stabilisers are exactly what we want to find. After choosing one subset $A$, its stabiliser $H$ acts on the elements of $A$. This second action is free: $hx=x$ forces $h=e$. Keeping the collection of subsets separate from the elements of one subset explains why the same stabiliser has both a lower and an upper bound on its order.

\begin{think}
Why would all Sylow $p$-subgroups be conjugate? Try $\mathfrak S_3$ and $p=2$: list its subgroups of order $2$ and conjugate one of them. Then guess $n_2$ and $n_3$ for $\mathfrak S_4$.
\end{think}

\begin{theorem}[Sylow II and III]\label{thm:sylow23}
Let $P$ be a Sylow $p$-subgroup of $G$.
\begin{enumerate}[label=(\alph*)]
\item Every $p$-subgroup of $G$ is contained in a conjugate of $P$.
\item The Sylow $p$-subgroups are exactly the conjugates of $P$.
\item $n_p=[G:N_G(P)]$, so $n_p$ divides $m$; and $n_p\equiv1\pmod p$.
\end{enumerate}
\end{theorem}
\begin{proof}
Write $|G|=p^a m$ with $p\nmid m$, so $|P|=p^a$.

\emph{(a) Every $p$-subgroup fits inside a conjugate of $P$.}
Let $Q\le G$ be any $p$-subgroup. It acts on the set of left cosets $G/P$ by
\[
q\cdot(gP)=(qg)P.
\]
This is well defined because a coset is a subset of $G$, and left multiplication sends it to another coset. The set has
\[
|G/P|=[G:P]=m
\]
elements. Orbit--stabiliser says that each orbit of the $p$-group $Q$ has a size dividing $|Q|$, hence a power of $p$. If every orbit had more than one point, then all orbit sizes would be divisible by $p$, and so would their sum $m$. This contradicts $p\nmid m$.

Consequently there is a coset $gP$ fixed by every element of $Q$. Being fixed means
\[
qgP=gP\quad(q\in Q).
\]
Multiplying the cosets on the left by $g^{-1}$ is equivalent to saying $g^{-1}qg\in P$ for every $q\in Q$. Thus
\[
g^{-1}Qg\subseteq P,\qquad\text{or}\qquad Q\subseteq gPg^{-1}.
\]

\emph{(b) All Sylow subgroups are conjugate.}
Every conjugate $gPg^{-1}$ has order $p^a$, so it is another Sylow $p$-subgroup. Conversely, let $Q$ itself be Sylow. Part (a) places it inside $gPg^{-1}$ for some $g$. Both groups have $p^a$ elements; an inclusion between two finite sets of the same size is equality. Hence $Q=gPg^{-1}$.

\emph{(c) Count the conjugates.}
The action of $G$ on its Sylow $p$-subgroups by conjugation is transitive by (b). The elements fixing $P$ form its normaliser
\[
N_G(P)=\{g\in G:gPg^{-1}=P\}.
\]
Orbit--stabiliser therefore gives
\[
n_p=[G:N_G(P)].
\]
Since $P\subseteq N_G(P)$, Lagrange gives
\[
[G:P]=[G:N_G(P)]\,[N_G(P):P],
\]
so $n_p$ divides $[G:P]=m$.

It remains to prove the congruence. Let $P$ act by conjugation on the same set of Sylow subgroups. Each orbit under $P$ has size a power of $p$. Modulo $p$, only orbits of size $1$ matter, so we must count the Sylow subgroups fixed by \emph{all} elements of $P$.

A Sylow subgroup $Q$ is fixed precisely when
\[
xQx^{-1}=Q\quad\text{for every }x\in P,
\]
that is, when $P\subseteq N_G(Q)$. In this situation both $P$ and $Q$ lie in $N_G(Q)$, and both have the largest possible $p$-power order $p^a$. They are therefore Sylow $p$-subgroups \emph{of the normaliser} too. Part (b), applied to the finite group $N_G(Q)$, says they are conjugate by some element of $N_G(Q)$. Yet $Q$ is normal in its normaliser, so all its conjugates there equal $Q$ itself. It follows that $P=Q$.

Thus $P$ is the unique fixed Sylow subgroup. All the other $P$-orbits have sizes divisible by $p$, and the one fixed orbit has size $1$. Adding their sizes gives
\[
n_p\equiv1\pmod p.
\]
\end{proof}

\paragraph{Why this works.}
The congruence counts subgroups, rather than their elements or their generators. Distinct subgroups of prime order intersect only in the identity, which is why the opening count for order $30$ works. Distinct subgroups of a larger prime-power order can overlap: in $\Z/4\Z\times\Z/2\Z$, the subgroups generated by $(1,0)$ and $(1,1)$ both have order $4$ and share $(2,0)$. Before multiplying a Sylow count by the number of nonidentity elements in one subgroup, we must prove that those elements are not being counted twice.

\begin{corollary}\label{cor:sylownormal}
A Sylow $p$-subgroup $P$ is normal iff $n_p=1$. In that case $P$ contains every $p$-subgroup of $G$, and every element of order a power of $p$.
\end{corollary}
\begin{proof}
Sylow's conjugacy theorem says that every Sylow
$p$-subgroup is of the form $gPg^{-1}$ for some $g\in G$.

If $P$ is normal, then $gPg^{-1}=P$ for all $g$.
So all Sylow $p$-subgroups coincide and $n_p=1$.
Conversely, if $n_p=1$, every conjugate $gPg^{-1}$
is a Sylow subgroup and must equal $P$.
This is exactly the definition of normality.

Suppose these equivalent conditions hold.
The Sylow containment theorem says that any
$p$-subgroup $H$ lies in a conjugate of $P$.
There is only one such conjugate, namely $P$,
so $H\subseteq P$.

If $x$ has order $p^j$, then
$\langle x\rangle$ is a subgroup of
order $p^j$, so it lies in $P$.
In particular $x\in P$.
This also covers $j=0$, the identity element.
\end{proof}

\begin{corollary}\label{cor:pk}
If $k\ge0$ is an integer and $p^k$ divides $|G|$, then $G$ has a subgroup of order $p^k$.
\end{corollary}
\begin{proof}
Write $|G|=p^am$ with $p\nmid m$ and $a\ge k$.
Sylow's first theorem gives a subgroup $P$
of order $p^a$. We will show by induction
on $a$ that \emph{every} group of order $p^a$
has subgroups of orders $1,p,\ldots,p^a$.

The case $a=0$ is clear, since the trivial
group contains its subgroup of order $1$.
For $a\ge1$, the centre $Z(P)$ is non-trivial
by Corollary~\ref{cor:pgroup}. Its order
divides $p^a$, so it is itself divisible
by $p$. By Cauchy's theorem it contains
an element $z$ of order $p$.
Let $C=\langle z\rangle$. Since $C$ is
central, it is normal in $P$.

The quotient $P/C$ has order $p^{a-1}$.
For $1\le j\le a$, the induction hypothesis
gives a subgroup $\bar K\le P/C$ of order
$p^{j-1}$. The preimage $K$ of $\bar K$
under the quotient map is a subgroup of $P$
containing $C$. Its size is
\[
|K|=|C|\,|\bar K|=p\cdot p^{j-1}=p^j,
\]
because every fibre of the quotient map
has exactly $|C|$ elements.
For $j=0$, take $\{e\}$.
This finishes the induction.

Set $j=k$. The resulting subgroup of $P$
is also a subgroup of $G$, with order $p^k$.
This does not claim that every divisor of
$|G|$ occurs as a subgroup order.
\end{proof}

The theorems restrict a Sylow count in two ways at once: it must divide an index and equal $1$ modulo the chosen prime. Even a small group rarely leaves many possibilities. Let us see what those arithmetic restrictions force.

\section{Using Sylow}

The method: factor $|G|$, list for each $p$ the divisors of $m$ that are $\equiv1\pmod p$, and count elements.

\begin{proposition}\label{prop:thirty}
No group of order $30$ is simple.
\end{proposition}
\begin{proof}
Suppose $|G|=30=2\cdot3\cdot5$.
Let $n_5$ and $n_3$ denote the numbers
of Sylow $5$- and $3$-subgroups.
Sylow's congruences and divisibility conditions give
\[
n_5\mid6,\quad n_5\equiv1\pmod5
\qquad\Longrightarrow\qquad n_5\in\{1,6\},
\]
and
\[
n_3\mid10,\quad n_3\equiv1\pmod3
\qquad\Longrightarrow\qquad n_3\in\{1,10\}.
\]

Assume neither number equals one. Then there
are six different subgroups of order five.
Two different such subgroups intersect only
at the identity: the intersection has order dividing five.
If its order were five, it would equal both subgroups,
contrary to their being different. Each contributes four distinct
elements of order five, so $G$ has
\[
6(5-1)=24
\]
different elements of order five.

The same reasoning for the ten distinct
subgroups of order three supplies
\[
10(3-1)=20
\]
different elements of order three.
An element cannot have both orders,
so these sets are disjoint.
We have counted $24+20=44$ non-identity
elements in a group with only $30$
elements, an impossibility.

Therefore $n_5=1$ or $n_3=1$.
By Corollary~\ref{cor:sylownormal},
the corresponding unique Sylow subgroup
is normal. Its order is $5$ or $3$,
so it is non-trivial and proper.
Thus $G$ is not simple.
\end{proof}

\begin{theorem}[groups of order $pq$]\label{thm:pq}
Let $p<q$ be primes and $|G|=pq$.
\begin{enumerate}[label=(\alph*)]
\item If $p\nmid q-1$, then $G\cong\Z/pq\Z$.
\item If $p\mid q-1$, then $G$ is either cyclic or isomorphic to the non-abelian semidirect product $\Z/q\Z\rtimes_u\Z/p\Z$, where $b\in\Z/p\Z$ acts on $\Z/q\Z$ as multiplication by $u^b$ and $u$ is any element of order $p$ in $(\Z/q\Z)^*$; this group does not depend on the choice of $u$.
\end{enumerate}
In particular every group of order $2q$ ($q$ an odd prime) is $\Z/2q\Z$ or the dihedral group $D_q$.
\end{theorem}
\begin{proof}
Let $P$ and $Q$ be Sylow subgroups of orders
$p$ and $q$ respectively. Because $q>p$,
the number $n_q$ divides $p$ and satisfies
$n_q\equiv1\pmod q$. Its only possible
value is $1$: the alternative $p$ is too
small to be congruent to $1$ modulo $q$.
Thus $Q$ is normal in $G$.

The groups $P$ and $Q$ have prime orders, hence
are cyclic. Write $P=\langle b\rangle$ and
$Q=\langle a\rangle$. Their intersection is
trivial, because its order divides both
distinct primes. The products $xy$ with
$x\in Q$, $y\in P$ are all distinct:
$xy=x'y'$ implies $x'^{-1}x=y'y^{-1}$
lies in $Q\cap P=\{e\}$.
There are $pq$ products, so $QP=G$.

Proposition~\ref{prop:semidirect} therefore
identifies $G$ with a semidirect product
$Q\rtimes P$. Conjugation by $b$
is an automorphism of the cyclic group $Q$,
so it sends its generator $a$ to $a^u$,
where $u$ is an invertible residue modulo $q$.
Since $b^p=e$, applying this automorphism
$p$ times gives
\[
u^p\equiv1\pmod q.
\]
The group $(\mathbb Z/q\mathbb Z)^*$ of such
residues is cyclic of order $q-1$ (Volume~2, Theorem~5.3).

\emph{(a) No non-trivial action when
$p\nmid(q-1)$.}
The order of $u$ divides both $p$ and
$q-1$. If $p\nmid(q-1)$ it must be $1$,
so $u=1$. Thus $b$ commutes with $a$,
and consequently $P$ commutes with $Q$.
The semidirect product is a direct product:
\[
G\cong\mathbb Z/q\mathbb Z\times
\mathbb Z/p\mathbb Z.
\]
The factors have coprime orders, and the
Chinese remainder theorem gives
\[
\boxed{G\cong\mathbb Z/pq\mathbb Z}.
\]

\emph{(b) Actions when $p\mid(q-1)$.}
Now the cyclic group $(\mathbb Z/q\mathbb Z)^*$
contains exactly one subgroup of order $p$.
An element satisfying $u^p=1$ either
is the identity or has order $p$.
The first case gives the direct product as
above. In the second case, choose $u$ of
order $p$ and define the action
\[
b\,a\,b^{-1}=a^u.
\]
The resulting semidirect product exists
by Proposition~\ref{prop:semidirect}(a).
It is not abelian, because $a^u\ne a$.

All non-trivial actions give isomorphic groups.
Indeed, any two elements of order $p$
in the cyclic automorphism group generate
the same subgroup. They are of the form
$u$ and $u^k$ with $\gcd(k,p)=1$.
Replacing the generator $b$ of $P$
by $b^k$ changes the displayed action
to exponent $u^k$ without changing $G$.

Hence exactly two isomorphism types occur
when $p\mid q-1$: the cyclic group and
the unique non-abelian semidirect product.
For $p=2$, $u=-1$ is the only element
of order two in $(\mathbb Z/q\mathbb Z)^*$,
and the non-abelian case is the dihedral
group $D_q$ of order $2q$.
\end{proof}

So a group of order $15$, $33$, $35$, $51$ or $65$ is cyclic, while order $21$ has exactly two groups, $\Z/21\Z$ and $\Z/7\Z\rtimes\Z/3\Z$ . Together with Exercise~\ref{ex:p2} in Chapter~3, Exercise~\ref{ex:order8} in Chapter~2, and Exercise~\ref{ex:order12} below, this gives the complete list of groups of order at most $15$, in Figure~\ref{fig:small}. For order $6$ the non-abelian group is $D_3$, and $D_3\cong\mathfrak S_3$.

\begin{figure}[ht]
\centering
\fig{ch05-small}
\caption{All groups of order $1$ to $15$ up to isomorphism, abelian ones in charcoal and non-abelian ones in coral ($\Z_n$ stands for $\Z/n\Z$). ``Prime'' is Corollary~\ref{cor:lagrange}, ``$p^2$'' is Exercise~\ref{ex:p2}, ``$2p$'' and ``$pq$'' are Theorem~\ref{thm:pq}, order $8$ is Exercise~\ref{ex:order8} and order $12$ is Exercise~\ref{ex:order12}. Order $16$ has $14$ groups, and order $2^{10}=1024$ has $49\,487\,367\,289$: the corrected enumeration of order $1024$ alone already dwarfs the groups in this small table (see Burrell in the sources).}
\label{fig:small}
\end{figure}

\begin{think}
For $|G|=12$: list the possible values of $n_2$ and $n_3$. Which combinations can occur together? Which group has $n_3=4$?
\end{think}

Sylow's theorems tell us when certain prime-power pieces exist. Another question will soon be more important: can a group be broken repeatedly through normal subgroups until every remaining quotient is abelian? That is the notion of solvability used in Galois theory.

\section{Solvable groups}

\begin{definition}
$G$ is \emph{solvable} if there is a chain of subgroups
\[
\{e\}=G_0\trianglelefteq G_1\trianglelefteq\dots\trianglelefteq G_r=G,
\]
each normal in the next (not necessarily in $G$), with every quotient $G_i/G_{i-1}$ abelian.
\end{definition}

The name comes from Galois theory: for a polynomial over a subfield of $\mathbb C$, Chapter~10 will prove that solvability by radicals is equivalent to solvability of its Galois group. This is an outlook; we use only the group definition here. Abelian groups are solvable ($r=1$), and so is $\mathfrak S_4$:
\[
\{e\}\trianglelefteq V\trianglelefteq A_4\trianglelefteq\mathfrak S_4,
\]
with $V$ the Klein group of Exercise~\ref{ex:v4} and quotients $V$, $A_4/V\cong\Z/3\Z$, $\mathfrak S_4/A_4\cong\Z/2\Z$, all abelian.

For $n\ge4$, $A_n$ is nonabelian. The two even permutations
$(1\,2\,3)$ and $(1\,2\,4)$ have products $(1\,3)(2\,4)$ and
$(1\,4)(2\,3)$ in the two orders, so they do not commute.

\begin{proposition}\label{prop:solv}
\begin{enumerate}[label=(\alph*)]
\item Subgroups and quotients of a solvable group are solvable.
\item If $N\trianglelefteq G$ and both $N$ and $G/N$ are solvable, then $G$ is solvable.
\item Every $p$-group is solvable. Every group of order $pq$, with $p,q$ prime, is solvable.
\item A non-abelian simple group is not solvable. In particular $A_n$ and $\mathfrak S_n$ are not solvable for $n\ge5$.
\end{enumerate}
\end{proposition}
\begin{proof}
A group is solvable if it has a chain
\[
\{e\}=G_0\trianglelefteq G_1
\trianglelefteq\cdots\trianglelefteq G_r=G
\]
with every quotient $G_i/G_{i-1}$ abelian.
We build the required chains using this definition.

\emph{(a) Subgroups and quotients.}
Let $H\le G$ and set $H_i=H\cap G_i$.
These groups form a chain from $\{e\}$ to $H$.
Because $G_{i-1}\trianglelefteq G_i$,
conjugation by any element of $H_i$
preserves $H_{i-1}$, so
$H_{i-1}\trianglelefteq H_i$.
The homomorphism
\[
H_i\to G_i/G_{i-1},
\qquad h\mapsto hG_{i-1},
\]
has kernel $H_{i-1}$. By the first
isomorphism theorem, $H_i/H_{i-1}$
is isomorphic to a subgroup of
the abelian group $G_i/G_{i-1}$,
hence is abelian. Thus $H$ is solvable.

Now let $N\trianglelefteq G$ and let
$\pi:G\to G/N$ be the quotient map.
Set $K_i=\pi(G_i)$. The $K_i$ form a
chain from $\{N\}$ to $G/N$.
Normality of $G_{i-1}$ in $G_i$
implies normality of $K_{i-1}$ in $K_i$.
Moreover the surjective map
$G_i\to K_i/K_{i-1}$ kills $G_{i-1}$,
so $K_i/K_{i-1}$ is a homomorphic
image of the abelian group $G_i/G_{i-1}$.
It is therefore abelian. Thus $G/N$
is solvable.

\emph{(b) Solvable normal subgroup and quotient.}
Assume $N\trianglelefteq G$, and both
$N$ and $G/N$ are solvable. Choose a
solvable chain
\[
\{e\}=N_0\trianglelefteq\cdots
\trianglelefteq N_s=N
\]
inside $N$, and another
\[
\{N\}=K_0\trianglelefteq\cdots
\trianglelefteq K_t=G/N.
\]
Write $L_i=\pi^{-1}(K_i)$.
Then $L_0=N$, $L_t=G$, and
$L_{i-1}\trianglelefteq L_i$.
The quotient map identifies
\[
L_i/L_{i-1}\cong K_i/K_{i-1},
\]
which is abelian by hypothesis.
Joining the two chains at $N=L_0$
therefore gives a solvable chain for $G$.

\emph{(c) The prime-power and $pq$ cases.}
Induct on the order of a $p$-group $P$.
The trivial group is solvable. A non-trivial
$p$-group has non-trivial centre $Z(P)$
by Corollary~\ref{cor:pgroup}.
The centre is abelian, hence solvable,
and $P/Z(P)$ is a smaller $p$-group,
solvable by induction. Apply (b) to
conclude that $P$ is solvable.

If the two primes coincide, the order is $p^2$, already covered.
Otherwise interchange their names if necessary so that $p<q$.
For a group of order $pq$ with $p<q$,
Theorem~\ref{thm:pq} supplies a normal
cyclic subgroup $Q$ of order $q$.
Its quotient has order $p$, so it is
also cyclic. Both are abelian and hence
solvable. Again use (b).

\emph{(d) Simple groups.}
Suppose $G$ is simple and non-abelian
but admits a solvable chain. Remove
any consecutive repeated terms.
The penultimate group $G_{r-1}$ is
a proper normal subgroup of $G$.
Simplicity forces $G_{r-1}=\{e\}$.
Then the final quotient
$G/G_{r-1}\cong G$ is abelian,
contrary to the assumption.

Thus a non-abelian simple group
cannot be solvable. In particular,
$A_n$ is simple and non-abelian
for $n\ge5$, so it is not solvable.
As $A_n\le\mathfrak S_n$,
part (a) also excludes solvability
of $\mathfrak S_n$ for $n\ge5$.
\end{proof}

\begin{remark}
Two external theorems extend (c). Burnside (1904): a group of order $p^\alpha q^\beta$, for primes $p,q$ and nonnegative integers $\alpha,\beta$, is solvable. Its classical proof uses character theory, whose basics we develop in Chapter~11. Feit and Thompson (1963): every finite group of odd order is solvable; their paper occupies $255$ pages. We do not prove either theorem here. The finite argument below uses only the Sylow and solvability results already earned.
\end{remark}

\subsection{The least nonabelian simple group}\label{tool:smallest-simple}
Finite Sylow counts prove that $A_5$ has the least order of a nonabelian simple group.
First every group of order less than $60$ is solvable. Induct on its order
$N$. The trivial group and the prime-power groups are solvable by Proposition~\ref{prop:solv}. Suppose $N<60$ is not a prime power, and every group of smaller positive order is solvable.

A Sylow count has a two-line shape. Write $N=p^am$ with $p\nmid m$. Then $n_p$ divides $m$ and $n_p\equiv1\pmod p$. If the only such positive integer is $1$, the Sylow $p$-subgroup is unique, hence normal. It is a solvable prime-power group, the quotient has smaller order and is solvable by induction, and Proposition~\ref{prop:solv}(b) makes $G$ solvable.

That test covers three unlike rows, and then the rest of the table below.
For order $6$ and the prime $3$, the cofactor is $2$, and $1+3=4$ does not divide $2$, so $n_3=1$.
For order $40$ and the prime $5$, the cofactor is $8$. The next candidate is $6$, and $6$ does not divide $8$, so $n_5=1$.
For order $45$ and the prime $5$, the cofactor is $9$, and again $6$ does not divide $9$, so $n_5=1$.
For every other row, the cofactor $m=N/p^a$ is strictly less than $p+1$. The only positive divisor of $m$ that is $1$ modulo $p$ is therefore $1$.

The same arithmetic forces a normal Sylow subgroup for the listed prime at each of the following orders:
\begin{center}\small
\begin{tabular}{c|l}
prime & orders $N$\\\hline
$3$ & $6,18,54$\\
$5$ & $10,15,20,40,45,50$\\
$7$ & $14,21,28,35,42$\\
$11$ & $22,33,44,55$\\
$13$ & $26,39,52$\\
$17$ & $34,51$\\
$19$ & $38,57$\\
$23$ & $46$\\
$29$ & $58$
\end{tabular}\end{center}
Each listed row is the cofactor test above, so the Sylow subgroup is normal.
It is a solvable prime-power group, and
its quotient is solvable by induction; Proposition~\ref{prop:solv}(b) applies.
The remaining non-prime-power orders below $60$ are
$12,24,30,36,48,56$.

For $12,24,36$, the number of Sylow $3$-subgroups is $1$ or $4$.
The first case gives a normal subgroup as before. In the second, conjugation
on the four Sylow subgroups is a homomorphism $\varphi:G\to\mathfrak S_4$.
The order of the image divides $|\mathfrak S_4|=24$. The action is nontrivial,
so the kernel $K=\ker\varphi$ is a proper normal subgroup and $|K|<|G|$.
The trivial group is solvable, and if $K$ is nontrivial then induction applies, so $K$ is solvable. The image is solvable, as a subgroup of the solvable group $\mathfrak S_4$.
The first isomorphism theorem identifies $G/K$ with that image, and the extension result makes $G$ solvable.

For order $30$, $n_5$ divides $6$ and is $1$ modulo $5$, so it is $1$ or $6$.
Also $n_3$ divides $10$ and is $1$ modulo $3$. The integer $4$ does not divide $10$, so $n_3$ is $1$ or $10$.
If neither subgroup is normal, the counts are $6$ and $10$.
They contribute $6\cdot4=24$ and $10\cdot2=20$ distinct nonidentity elements, which cannot fit in a group of order $30$.
Thus one is normal and induction applies.

For order $48$, $n_3$ is $1$, $4$ or $16$. The first two cases are the argument just given for order $36$: a normal Sylow subgroup, or a homomorphism into $\mathfrak S_4$ whose image order divides $24$ and whose kernel is proper.
In the last case, the $16$ subgroups of order $3$ contain $16\cdot2=32$ distinct nonidentity elements.
Only $16$ elements remain, including the identity.
Every Sylow $2$-subgroup has order $16$ and contains none of those elements of order $3$, so it is exactly the remaining set. It is unique and normal.
For order $56$, $n_7$ is $1$ or $8$. In the latter case its $8\cdot6=48$ nonidentity elements leave exactly the $8$ elements of a unique normal Sylow $2$-subgroup.
Induction and the extension result complete all orders below $60$.
As a nonabelian simple group cannot be solvable, $A_5$ has the least possible order of such a group.

There is also only one simple group of order $60$ up to isomorphism.
A nontrivial homomorphism from a simple group is injective: its kernel is a normal subgroup, and the homomorphism is not the trivial one.
Let $G$ be simple of order $60$. Then $n_5$ divides $12$ and is $1$ modulo $5$, so it is $1$ or $6$. Simplicity forces $n_5=6$.
The candidates for $n_3$ are the divisors of $20$ that are $1$ modulo $3$, namely $1$, $4$ and $10$.
The count $1$ would be normal. The count $4$ would be a nontrivial, hence injective, homomorphism into $\mathfrak S_4$, and $60$ does not embed into a group of order $24$. Thus $n_3=10$.
The candidates for $n_2$ are the divisors of $15$, namely $1$, $3$, $5$ and $15$.
The count $1$ is normal. The count $3$ would be an injective homomorphism into $\mathfrak S_3$, and $60$ does not divide $6$. Thus $n_2$ is $5$ or $15$.

If $n_2=5$, conjugation on these five subgroups is a nontrivial homomorphism $G\to\mathfrak S_5$, hence injective.
The image has index $2$. In a quotient of order $2$, the image of every square is the identity, so the subgroup contains every permutation square.
A $3$-cycle $\tau$ satisfies $\tau^3=e$ and $\tau\ne e$, so $\tau^2\ne e$ and $(\tau^2)^2=\tau^4=\tau$.
Thus every $3$-cycle is a square. Theorem~\ref{thm:an} says the $3$-cycles generate $A_5$, so the image contains $A_5$.
The two groups both have order $60$, and the image is exactly $A_5$.

Suppose instead $n_2=15$. The elements of order $5$ number $6\cdot4=24$, and the elements of order $3$ number $10\cdot2=20$.
Together with the identity that is $45$ elements, leaving at most $15$ others.
If a Sylow $2$-subgroup were cyclic of order $4$, Sylow conjugacy would make all $15$ of them cyclic.
Each has $\varphi(4)=2$ generators of order $4$, and those generators are disjoint from one cyclic group to the next, giving $15\cdot2=30$ elements of order $4$.
That is impossible. Hence every Sylow $2$-subgroup is a Klein four-group.
Each has three involutions, so the $15$ subgroups give $15\cdot3=45$ incidences with involutions, among at most $15$ elements.
Some involution $a$ therefore lies in more than one Sylow subgroup.
Chapter~1 classifies the groups of order $4$: each is abelian. Any Sylow subgroup containing $a$ therefore lies in $C_G(a)$.
Their union has more than four elements, so $|C_G(a)|>4$. It is a multiple of $4$ dividing $60$, hence $12$, $20$ or $60$.
Order $60$ would make $a$ central, contradicting simplicity.
Order $20$ would give a nontrivial, hence injective, action on three cosets, and $60$ does not embed in $\mathfrak S_3$.
Thus $|C_G(a)|=12$. Acting on its five cosets is again an injective homomorphism to $\mathfrak S_5$, and the index-$2$ argument above gives $G\cong A_5$.
This proves uniqueness as well as minimality.

\begin{beauty}
The group-theoretic distinction behind the quartic and quintic is now proved: every subgroup of $\mathfrak S_4$ is solvable, whereas $\mathfrak S_5$ is not. Chapter~10 will show that the polynomial $x^5-4x+2$ has Galois group $\mathfrak S_5$ and connect these group properties with formulas by radicals over subfields of $\mathbb C$.

There is also a computational use for prime-power pieces. If $g$ generates a finite cyclic group of order $N$ and $h=g^x$, recovering $x$ modulo $N$ is a \emph{discrete logarithm}. Suppose $N>1$ is factored. For each prime power $\ell^a$ in its factorisation, raise both sides to $N/\ell^a$:
\[
h^{N/\ell^a}=\bigl(g^{N/\ell^a}\bigr)^x.
\]
The new base has order $\ell^a$, so its discrete logarithm determines $x$ modulo $\ell^a$. The Chinese remainder theorem then combines those answers into $x$ modulo $N$. For $N=1$ the only answer is $0$ modulo $1$.
This is the prime-power reduction in the Pohlig--Hellman algorithm; the efficient work inside each prime-power subgroup goes beyond this chapter. Its cryptographic significance is that a large group order alone is insufficient when all its prime factors are small.
\end{beauty}

\section{Exercises}

\begin{exercise}[$\star$]\label{ex:sylowS4}
Find $n_2$ and $n_3$ for $\mathfrak S_4$, and describe its Sylow subgroups.
\end{exercise}

\begin{exercise}[$\star$]\label{ex:order45}
Show that every group of order $45$ is abelian. Which abelian groups of order $45$ are there?
\end{exercise}

\begin{exercise}[$\star$]\label{ex:nosimple}
Show that there is no simple group of order (a) $56$; (b) $200$; (c) $36$. For (c), let $G$ act on its Sylow $3$-subgroups.
\end{exercise}

\begin{exercise}[$\star\star$]\label{ex:S5sub}
Show that $\mathfrak S_5$ has no subgroup of order $15$, $30$ or $40$, although these numbers divide $120$.
\end{exercise}

\begin{exercise}[$\star\star$]\label{ex:normalizer}
Let $G$ be finite, $P$ a Sylow $p$-subgroup of $G$ and $M=N_G(P)$. Show that $N_G(M)=M$.
\end{exercise}

\begin{exercise}[$\star\star$]\label{ex:derived}
With the derived subgroup of Exercise~\ref{ex:commutator}, define $G^{(0)}=G$ and $G^{(i+1)}=\bigl(G^{(i)}\bigr)'$.
\begin{enumerate}[label=(\alph*)]
\item Show that $G$ is solvable iff $G^{(k)}=\{e\}$ for some $k$.
\item Compute the derived series of $\mathfrak S_4$ and of $\mathfrak S_5$.
\end{enumerate}
\end{exercise}

\begin{exercise}[$\star\star$]\label{ex:sylowGL2}
Let $p$ be prime and $G=GL_2(\Z/p\Z)$, of order $(p^2-1)(p^2-p)$ (Exercise~\ref{ex:gl2}). Show that the matrices $\begin{psmallmatrix}1&x\\0&1\end{psmallmatrix}$ form a Sylow $p$-subgroup $U$, that $N_G(U)$ is the group $B$ of invertible upper triangular matrices, and that $n_p=p+1$. Interpret the $p+1$ Sylow subgroups geometrically.
\end{exercise}

\begin{exercise}[$\star\star\star$]\label{ex:order12}
Show that there are exactly five groups of order $12$ up to isomorphism: $\Z/12\Z$, $\Z/2\Z\times\Z/6\Z$, $A_4$, $D_6$ and $\Z/3\Z\rtimes\Z/4\Z$ (where a generator of $\Z/4\Z$ acts on $\Z/3\Z$ by $x\mapsto-x$).
\end{exercise}

\section{Hints}

\begin{hint}{ex:sylowS4}
$n_3$: count the $3$-cycles. $n_2$: use the action of $\mathfrak S_4$ on the three splittings into pairs (Exercise~\ref{ex:v4}).
\end{hint}
\begin{hint}{ex:order45}
Both Sylow subgroups are normal. Then use Exercise~\ref{ex:p2}.
\end{hint}
\begin{hint}{ex:nosimple}
(a) If $n_7=8$, count the elements of order $7$; what is left? (b) $n_5$ divides $8$. (c) $n_3\in\{1,4\}$; a homomorphism $G\to\mathfrak S_4$ cannot be injective.
\end{hint}
\begin{hint}{ex:S5sub}
For a subgroup $H$ of index $k$, let $\mathfrak S_5$ act on the $k$ cosets of $H$; the kernel is a normal subgroup contained in $H$. Use Exercise~\ref{ex:normalSn}. For $15$, use Theorem~\ref{thm:pq} and Exercise~\ref{ex:landau}.
\end{hint}
\begin{hint}{ex:normalizer}
If $g$ normalises $M$, then $gPg^{-1}$ is a Sylow $p$-subgroup of $M$. How many does $M$ have?
\end{hint}
\begin{hint}{ex:derived}
(a) If $G_r/G_{r-1}$ is abelian, $G'\subseteq G_{r-1}$; iterate. (b) The commutator of $(a\,b)$ and $(a\,c)$ is a $3$-cycle; in $A_4$, the quotient by $V$ is abelian.
\end{hint}
\begin{hint}{ex:sylowGL2}
The vectors fixed by every element of $U$ form the line spanned by $e_1$. If $gUg^{-1}=U$, where does $g$ send this line?
\end{hint}
\begin{hint}{ex:order12}
If $n_3=4$, let $G$ act on its four Sylow $3$-subgroups. If $n_3=1$, $G$ is a semidirect product $\Z/3\Z\rtimes Q$ with $|Q|=4$, and $\operatorname{Aut}(\Z/3\Z)$ has order $2$.
\end{hint}

\section{Solutions}

\begin{solution}{ex:sylowS4}
\idea Factor $|\mathfrak S_4|$, and treat the two primes separately. Three-cycles determine the Sylow $3$-subgroups. For the Sylow $2$-subgroups, use the action on the three partitions into pairs from Exercise~\ref{ex:v4}.

\writeup Since
\[
|\mathfrak S_4|=4!=24=2^3\cdot3,
\]
a Sylow $3$-subgroup has order $3$ and a Sylow $2$-subgroup has order $8$.

\emph{Sylow $3$-subgroups.} Every subgroup of order $3$ is cyclic and consists of $e$, a $3$-cycle, and its inverse. There are
\[
\binom43\cdot(3-1)!=4\cdot2=8
\]
three-cycles in $\mathfrak S_4$. Each subgroup of order $3$ contains exactly two of them, and distinct such subgroups meet only in the identity. Consequently
\[
\boxed{n_3=8/2=4}.
\]
This satisfies Sylow III: $4\mid8$ and $4\equiv1\pmod3$.

\emph{Sylow $2$-subgroups.} Let $\mathfrak S_4$ act on the three partitions $P_1,P_2,P_3$ of four letters into two unordered pairs, as in Exercise~\ref{ex:v4}. This action is transitive. Orbit--stabiliser gives, for each partition $P_i$,
\[
|\operatorname{Stab}(P_i)|=\frac{24}{3}=8.
\]
So each stabiliser is a Sylow $2$-subgroup.

The action map $\mathfrak S_4\to\mathfrak S_3$ is onto by Exercise~\ref{ex:v4}. Each stabiliser is the preimage of a point stabiliser in $\mathfrak S_3$; those three order-two subgroups are distinct, so their preimages are distinct too. Conjugation carries the stabiliser of one partition to that of another, so the three groups are conjugate. Sylow III also gives $n_2\mid3$; since there are at least three, there are exactly three:
\[
\boxed{n_2=3}.
\]

Each of these order-eight stabilisers is isomorphic to $D_4$, the symmetry group of a square. To see this geometrically, put the four letters at the vertices of a square so that the two pairs in $P_i$ form its two diagonals. A permutation preserving the set of diagonals is exactly a symmetry of that square. Thus its stabiliser has the eight dihedral symmetries, consisting of four rotations and four reflections.

We have described all the Sylow subgroups: four cyclic groups of order $3$ and three dihedral groups of order $8$.
\end{solution}

\begin{solution}{ex:order45}
\idea The divisibility and congruence conditions in Sylow III are unusually restrictive when the group order is $3^2\cdot5$. If both Sylow subgroups are unique, normality and their coprime orders let us reconstruct the whole group as a direct product.

\writeup Let $|G|=45=9\cdot5$. If $n_5$ is the number of Sylow $5$-subgroups, then
\[
n_5\mid9,\qquad n_5\equiv1\pmod5.
\]
Of the divisors $1,3,9$ of $9$, only $1$ is $1$ modulo $5$. Thus $n_5=1$, and the Sylow $5$-subgroup $P_5$ is normal.

Similarly,
\[
n_3\mid5,\qquad n_3\equiv1\pmod3.
\]
The divisors of $5$ are $1$ and $5$, and only $1$ is $1$ modulo $3$, so $n_3=1$. The Sylow $3$-subgroup $P_3$ of order $9$ is normal too.

Their intersection is a subgroup of both, so its order divides both $9$ and $5$, forcing $P_3\cap P_5=\{e\}$. The $9\cdot5$ products $xy$ with $x\in P_3$, $y\in P_5$ are distinct: $xy=x'y'$ forces $x'^{-1}x=y'y^{-1}$ to lie in the trivial intersection. Thus $P_3P_5$ has $45$ elements and equals $G$. Two normal subgroups with trivial intersection commute elementwise: for $x\in P_3$ and $y\in P_5$, the commutator $xyx^{-1}y^{-1}$ belongs to both $P_3$ and $P_5$, hence is $e$. The multiplication map
\[
P_3\times P_5\longrightarrow G,\quad(x,y)\mapsto xy
\]
is consequently an isomorphism.

By Exercise~\ref{ex:p2}, every group of order $9$ is abelian and is isomorphic either to $\Z/9\Z$ or to $(\Z/3\Z)^2$. A group of order $5$ is cyclic. Hence there are exactly two possibilities:
\[
G\cong\Z/9\Z\times\Z/5\Z\cong\boxed{\Z/45\Z}
\]
or
\[
G\cong(\Z/3\Z)^2\times\Z/5\Z
\cong\boxed{\Z/3\Z\times\Z/15\Z}.
\]
They are non-isomorphic: the cyclic group has an element of order $9$, whereas the other group has no element of order $9$. Both are abelian, proving the claim.
\end{solution}

\begin{solution}{ex:nosimple}
\idea Sylow's theorems tell us how many subgroups of the largest prime-power orders are possible. A unique Sylow subgroup is normal. When uniqueness fails, count the elements it forces or let the group permute those subgroups.

\writeup Recall that a simple group has no normal subgroups other than $\{e\}$ and itself. Every Sylow subgroup considered below has order strictly between $1$ and $|G|$, so finding a unique one will rule out simplicity.

\emph{(a) Order $56=2^3\cdot7$.} The number of Sylow $7$-subgroups must divide $8$ and be congruent to $1$ modulo $7$. Hence $n_7=1$ or $n_7=8$. The first case supplies a normal subgroup of order $7$.

In the second case, two distinct subgroups of order $7$ intersect only at the identity, since their intersection has order dividing $7$. They therefore account for $8(7-1)=48$ distinct non-identity elements. Exactly eight elements of $G$ remain outside this set, including the identity. Any Sylow $2$-subgroup has eight elements, none of which has order $7$, so it is precisely the set of eight remaining elements. It is therefore the \emph{only} Sylow $2$-subgroup, hence normal.

\emph{(b) Order $200=2^3\cdot5^2$.} A Sylow $5$-subgroup has order $25$. Its count divides $8$ and is $1$ modulo $5$. Of $1,2,4,8$, only $1$ has that property. The unique Sylow $5$-subgroup is normal and proper.

\emph{(c) Order $36=2^2\cdot3^2$.} The number $n_3$ divides $4$ and is congruent to $1$ modulo $3$, hence is $1$ or $4$. Again $n_3=1$ gives a proper normal subgroup.

Suppose $n_3=4$. Conjugation permutes the four Sylow $3$-subgroups, yielding a homomorphism
\[
\varphi:G\longrightarrow\mathfrak S_4.
\]
The action is transitive by Sylow II, so $\varphi$ is not trivial; thus its kernel is a proper normal subgroup. That kernel cannot be the identity, because an injective map would make a group of order $36$ a subgroup of $\mathfrak S_4$, which has only $24$ elements. Therefore the kernel is non-trivial as well. This also rules out simplicity.

All three orders are impossible for a simple group, but note the distinction: cases (a) and (b) find normal Sylow subgroups; case (c) finds a normal \emph{kernel}.
\end{solution}

\begin{solution}{ex:S5sub}
\idea Lagrange's theorem makes the three subgroup orders possible divisors, not guaranteed subgroup sizes. For orders $30$ and $40$, the action on a small number of cosets gives a homomorphism into a symmetric group too small to contain the image. For order $15$, use the internal structure of a group of order $3\cdot5$ and the possible cycle types in $\mathfrak S_5$.

\writeup Let $G=\mathfrak S_5$, so $|G|=120$.

\emph{No subgroup of order $40$.} If $H\le G$ had order $40$, its index would be $3$. The action of $G$ on the three left cosets $G/H$ would give
\[
\varphi:G\longrightarrow\mathfrak S_3.
\]
Its kernel $K$ is normal in $G$, and every element of $K$ fixes the coset $H$, hence $K\subseteq H$. By the first isomorphism theorem,
\[
|G/K|=|\varphi(G)|\le|\mathfrak S_3|=6,
\]
so $|K|\ge120/6=20$; yet $|K|\le40$. Exercise~\ref{ex:normalSn} tells us that a normal subgroup of $\mathfrak S_5$ has order only $1$, $60$ or $120$. None lies between $20$ and $40$, a contradiction.

\emph{No subgroup of order $30$.} Such an $H$ would have index $4$. The same coset action would yield a kernel $K\trianglelefteq G$ contained in $H$ with
\[
|K|\ge120/|\mathfrak S_4|=120/24=5,\qquad |K|\le30.
\]
Again none of $1,60,120$ is possible.

\emph{No subgroup of order $15$.} Suppose $H$ had order $15=3\cdot5$. Sylow III applied \emph{inside $H$} gives
\[
n_5\mid3,\ n_5\equiv1\pmod5
\quad\Longrightarrow\quad n_5=1,
\]
and
\[
n_3\mid5,\ n_3\equiv1\pmod3
\quad\Longrightarrow\quad n_3=1.
\]
So the subgroups $P_5$ and $P_3$ are both normal, have trivial intersection, and have product all of $H$. Their orders are prime, hence they are cyclic, and the internal direct product theorem gives
\[
H\cong\Z/5\Z\times\Z/3\Z\cong\Z/15\Z.
\]
This would give an element of order $15$ in $\mathfrak S_5$.

But a permutation has order equal to the lcm of its disjoint cycle lengths. To have order divisible by $15=3\cdot5$, it must include a cycle length divisible by $5$ and one divisible by $3$. In five letters, a $5$-cycle uses all the letters and leaves no room for a $3$-cycle; no cycle length at most five is itself divisible by $15$. Therefore no permutation of five letters has order $15$, a contradiction.

Thus $\mathfrak S_5$ has no subgroup of order $15$, $30$ or $40$. Lagrange's theorem only provides a \emph{necessary} divisibility condition for subgroup orders; it never asserted the converse.
\end{solution}

\begin{solution}{ex:normalizer}
\idea The normaliser $M=N_G(P)$ is the largest subgroup of $G$ in which $P$ is normal. Any element normalising $M$ sends $P$ to another Sylow subgroup of $M$, but $M$ has only one of those.

\writeup Let $P$ be a Sylow $p$-subgroup of $G$, of order $p^a$, and put
\[
M=N_G(P)=\{g\in G:gPg^{-1}=P\}.
\]
By definition, $P\trianglelefteq M$. It is still a Sylow $p$-subgroup when considered inside $M$: no $p$-subgroup of $M\subseteq G$ can have order greater than $p^a$.

We now show $N_G(M)\subseteq M$. Take $g\in N_G(M)$. Thus $gMg^{-1}=M$. Since $P\subseteq M$, conjugation gives
\[
gPg^{-1}\subseteq gMg^{-1}=M.
\]
This conjugate has order $p^a$, so it is another Sylow $p$-subgroup of $M$.

Because $P$ is normal in $M$, all its conjugates \emph{inside $M$} are equal to $P$. Sylow II says every Sylow $p$-subgroup of $M$ is one of those conjugates. Therefore $P$ is the unique Sylow $p$-subgroup of $M$, and in particular
\[
gPg^{-1}=P.
\]
That means $g\in N_G(P)=M$.

Conversely, every subgroup is contained in its own normaliser: $mMm^{-1}=M$ for all $m\in M$. So $M\subseteq N_G(M)$, and combining the inclusions gives
\[
\boxed{N_G(N_G(P))=N_G(P)}.
\]
The statement relies on $P$ being Sylow. It is not generally true for the normaliser of an arbitrary subgroup.
For example, in $D_4=\langle r,s\rangle$, let $H=\langle s\rangle$.
The relations $r^4=s^2=e$ and $srs=r^{-1}$ give
$N_{D_4}(H)=\{e,r^2,s,r^2s\}$. This subgroup has index $2$, so is normal
in $D_4$; its own normaliser is therefore all of $D_4$, strictly larger.
\end{solution}

\begin{solution}{ex:derived}
\idea The derived subgroup $H'$ is the smallest normal subgroup of $H$ whose quotient is abelian. The derived series repeatedly removes everything obstructing commutativity. Compare it with a solvable chain and then calculate in $\mathfrak S_4$ and $\mathfrak S_5$.

\writeup Write $G^{(0)}=G$ and $G^{(i+1)}=(G^{(i)})'$. By Exercise~\ref{ex:commutator}, each derived subgroup is normal in its predecessor, and
\[
G^{(i)}/G^{(i+1)}
\]
is abelian. Thus if $G^{(k)}=\{e\}$, the chain $G=G^{(0)}\supseteq\cdots\supseteq G^{(k)}=\{e\}$ satisfies the definition of solvability.

Conversely suppose $G$ admits a chain
\[
G=H_0\trianglerighteq H_1\trianglerighteq\cdots
\trianglerighteq H_r=\{e\},
\qquad H_i/H_{i+1}\text{ abelian}.
\]
We prove by induction that $G^{(i)}\subseteq H_i$. The assertion is clear for $i=0$. If it holds for $i$, every commutator of elements in $G^{(i)}$ lies in $H_i'$, so
\[
G^{(i+1)}=(G^{(i)})'\subseteq H_i'\subseteq H_{i+1};
\]
the last inclusion holds because $H_i/H_{i+1}$ is abelian. Thus $G^{(r)}\subseteq H_r=\{e\}$. This proves the equivalence.

For $\mathfrak S_n$ ($n\ge3$), the sign map has abelian image, so $\mathfrak S_n'\subseteq A_n$. Conversely, the commutator of two transpositions sharing one letter is a $3$-cycle, and all $3$-cycles are obtained by relabelling. Since $A_n$ is generated by $3$-cycles,
\[
\mathfrak S_n'=A_n.
\]

In $A_4$, the Klein subgroup $V_4$ is normal with abelian quotient $A_4/V_4\cong\Z/3\Z$, so $A_4'\subseteq V_4$. For the reverse inclusion, direct multiplication gives
$[(1\,2\,3),(1\,2\,4)]=(1\,2)(3\,4)$.
Conjugating this element by $(1\,2\,3)$ and its square gives the other
two double transpositions. Since $A_4'$ is normal in $A_4$, it contains
all three, so $A_4'=V_4$. As $V_4$ is abelian, $V_4'=\{e\}$. Therefore
\[
\boxed{\mathfrak S_4\trianglerighteq A_4
\trianglerighteq V_4\trianglerighteq\{e\}}
\]
is its derived series and $\mathfrak S_4$ is solvable.

In $A_5$, the derived subgroup is normal and non-trivial, since $A_5$ is non-abelian. Simplicity of $A_5$ forces $A_5'=A_5$. So
\[
\mathfrak S_5^{(0)}=\mathfrak S_5,\quad
\mathfrak S_5^{(1)}=A_5,\quad
\mathfrak S_5^{(i)}=A_5\quad(i\ge1).
\]
The series never reaches $\{e\}$, and $\mathfrak S_5$ is not solvable.
A group $H$ with $H'=H$ is called \emph{perfect}: $A_5$ is perfect,
whereas the simple abelian group $\mathbb Z/3\mathbb Z$ is not.
Perfect groups need not be simple. In $A_5\times A_5$, commutators are
computed in each coordinate. Commutators of elements in the first factor
generate $A_5\times\{e\}$, and those in the second generate
$\{e\}\times A_5$. Thus the product is perfect, but either factor is a
proper non-trivial normal subgroup. At the other boundary, the trivial
group is both perfect and solvable. These properties are different from
simplicity, which excludes the trivial group.
\end{solution}

\begin{solution}{ex:sylowGL2}
\idea The non-trivial matrices in $U$ are shears. Every shear fixes exactly one common line pointwise. Conjugation moves this line, so finding the normaliser is the same as finding its stabiliser.

\writeup Set $K=\mathbb F_p$ and
\[
U=\left\{u_x=\begin{pmatrix}1&x\\0&1\end{pmatrix}:x\in K\right\}.
\]
Multiplication gives $u_xu_y=u_{x+y}$ and $u_x^{-1}=u_{-x}$. Hence $U$ is a subgroup isomorphic to $(K,+)$, of order $p$. Since
\[
|GL_2(K)|=(p^2-1)(p^2-p)
=p(p-1)^2(p+1)
\]
and neither $p-1$ nor $p+1$ is divisible by $p$, the highest $p$-power dividing $|G|$ is $p$. Thus $U$ is a Sylow $p$-subgroup.

Let $B$ denote the invertible upper triangular matrices. For
\[
g=\begin{pmatrix}a&b\\0&d\end{pmatrix}\in B
\]
with $a,d\ne0$, direct multiplication yields
\[
gu_xg^{-1}=u_{(a/d)x}\in U.
\]
Therefore $B\subseteq N_G(U)$.

For the reverse inclusion, look at the vectors fixed by \emph{all} elements of $U$. If $v=(v_1,v_2)^{\mathsf T}$, then
\[
u_xv=(v_1+xv_2,v_2)^{\mathsf T}.
\]
This equals $v$ for every $x$ exactly when $v_2=0$. Thus the common fixed space is the line $L=Ke_1$.

If $gUg^{-1}=U$, then $g$ sends the common fixed space of $U$ onto itself: for any $v$ fixed by $U$ and any $u\in U$,
\[
u(gv)=g(g^{-1}ug)v=gv.
\]
Hence $gL=L$, which means that $ge_1$ has second coordinate $0$. So $g$ is upper triangular and belongs to $B$. We conclude
\[
\boxed{N_G(U)=B}.
\]

The group $B$ has $(p-1)$ choices for each non-zero diagonal entry and $p$ choices for its upper-right entry, hence
\[
|B|=p(p-1)^2.
\]
Sylow II identifies the number of Sylow $p$-subgroups with the index of their normaliser:
\[
\boxed{n_p=[G:B]=p+1}.
\]

There are exactly $(p^2-1)/(p-1)=p+1$ lines through the origin in $K^2$. The conjugate $gUg^{-1}$ fixes the line $gL$ pointwise and consists of the shears directed along that line. The correspondence between Sylow subgroups and lines is bijective because the normaliser equals the line stabiliser.
\end{solution}

\begin{solution}{ex:order12}
\idea Sylow's theorem gives only two choices for the number of subgroups of order $3$. One choice makes $G$ act faithfully on four objects; the other leaves a normal subgroup of order $3$, and then we must decide how a group of order $4$ acts on it.

\writeup Let $|G|=12=3\cdot4$. Sylow III says that the number $n_3$ of Sylow $3$-subgroups divides $4$ and is congruent to $1$ modulo $3$. Hence
\[
\boxed{n_3=1\quad\text{or}\quad n_3=4.}
\]

\emph{First case: four subgroups of order $3$.}
Call them $P_1,\ldots,P_4$. Conjugation permutes these subgroups, giving a homomorphism
\[
\varphi:G\longrightarrow\mathfrak S_4.
\]
The normaliser of each $P_i$ has index $n_3=4$, so
\[
|N_G(P_i)|=12/4=3.
\]
Since $P_i\subseteq N_G(P_i)$ and both have order $3$, we have $N_G(P_i)=P_i$.
An element in $\ker\varphi$ normalises every $P_i$, so belongs to their intersection. Distinct subgroups of order $3$ intersect only in the identity. Thus $\ker\varphi=\{e\}$.

We have embedded $G$ as a subgroup $H$ of $\mathfrak S_4$ with $|H|=12$. Such a subgroup has index $2$, hence is normal. The quotient map $\mathfrak S_4\to\mathfrak S_4/H$ is a non-trivial homomorphism to a group of order $2$. Every transposition must have the same image, because transpositions are conjugate. They generate $\mathfrak S_4$, so their image cannot be trivial. Therefore this quotient map agrees with the sign map, and
\[
H=A_4,\qquad G\cong A_4.
\]

\emph{Second case: one subgroup of order $3$.}
Write $P=\langle a\rangle$. It is normal, since every conjugate is another Sylow $3$-subgroup. Choose a Sylow $2$-subgroup $Q$; it has order $4$. The intersection $P\cap Q$ is trivial because its order divides both $3$ and $4$, and $PQ$ has $3\cdot4=12$ elements. Therefore
\[
G\cong P\rtimes_\varphi Q,
\qquad
\varphi:Q\longrightarrow\operatorname{Aut}(P)\cong\mathbb Z/2\mathbb Z.
\]
The classification of groups of order $4$ (Exercise~\ref{ex:order4}) leaves $Q\cong\mathbb Z/4\mathbb Z$ or $Q\cong(\mathbb Z/2\mathbb Z)^2$.

There are three possibilities for the action.

\begin{itemize}
\item \emph{Trivial action.} The semidirect product is a direct product. If $Q$ is cyclic, we obtain $\mathbb Z/12\mathbb Z$. If $Q$ is the Klein four-group, we obtain $\mathbb Z/2\mathbb Z\times\mathbb Z/6\mathbb Z$.

\item \emph{Non-trivial action with $Q$ cyclic.} A generator of $\mathbb Z/4\mathbb Z$ must send $a$ to $a^{-1}$, since inversion is the only non-identity automorphism of $P$. We obtain exactly one action and hence the group $\mathbb Z/3\mathbb Z\rtimes\mathbb Z/4\mathbb Z$ specified in the question.

\item \emph{Non-trivial action with $Q$ the Klein four-group.} The homomorphism $\varphi$ has a kernel of order $2$. Every subgroup of order $2$ in $Q$ can be carried to any other by an automorphism of $Q$ (permute its three non-identity elements). Consequently all three non-trivial actions give isomorphic semidirect products: if $\varphi'=\varphi\circ\alpha$ then $(p,q)\mapsto(p,\alpha^{-1}(q))$ is an isomorphism from $P\rtimes_\varphi Q$ to $P\rtimes_{\varphi'}Q$: in the target, $\varphi'(\alpha^{-1}(q))=\varphi(q)$, so the multiplication laws agree. One concrete model is $D_6=\langle r,s:r^6=s^2=e,\ srs=r^{-1}\rangle$: take $P=\langle r^2\rangle$ and $Q=\{e,r^3,s,r^3s\}$.
\end{itemize}

That is one group in the first case and four in the second.

\emph{Why these are five different groups.}
The two direct products are abelian; only $\mathbb Z/12\mathbb Z$ has an element of order $12$. The remaining three are non-abelian. The group $A_4$ has four Sylow $3$-subgroups, while the other two have one. Finally, $D_6$ has a non-cyclic Sylow $2$-subgroup, whereas $\mathbb Z/3\mathbb Z\rtimes\mathbb Z/4\mathbb Z$ has a cyclic Sylow $2$-subgroup. Thus none of the five is isomorphic to another.

\tools Sylow II and III; conjugation actions (Chapter~3); the sign map (Chapter~4); semidirect products (Chapter~2).

\leads This classification is a good warning: even at order $12$, knowing the prime factorisation of $|G|$ is not enough. We also need to understand how the prime-power subgroups fit together.
\end{solution}

\section*{Sources and further reading}

\begin{itemize}
\item L.~Sylow, \emph{Th\'eor\`emes sur les groupes de substitutions}, Math.\ Ann.\ 5 (1872), 584--594. The original theorems, stated for permutation groups.
\item H.~Wielandt, \emph{Ein Beweis f\"ur die Existenz der Sylowgruppen}, Arch.\ Math.\ 10 (1959), 401--402. The counting proof of Theorem~\ref{thm:sylow1}.
\item W.~Feit and J.~G.~Thompson, \emph{Solvability of groups of odd order}, Pacific J.\ Math.\ 13 (1963), 775--1029. The odd order theorem of the remark; formally verified in the Coq proof assistant in 2012.
\item H.~U.~Besche, B.~Eick and E.~A.~O'Brien, \emph{A millennium project: constructing small groups}, Internat.\ J.\ Algebra Comput.\ 12 (2002), 623--644. The small-group enumeration project; its reported count for order $1024$ was later corrected by Burrell.
\item D.~Burrell, \emph{On the number of groups of order 1024}, Comm.\ Algebra 50 (2022), 2408--2410, \url{https://doi.org/10.1080/00927872.2021.2006680}. Published online in 2021; the corrected count used in Figure~\ref{fig:small}.
\item S.~C.~Pohlig and M.~E.~Hellman, \emph{An improved algorithm for computing logarithms over $GF(p)$ and its cryptographic significance}, IEEE Trans.\ Inform.\ Theory 24 (1978), 106--110. The discrete-logarithm algorithm.
\item D.~S.~Dummit and R.~M.~Foote, \emph{Abstract Algebra}, 3rd ed., Wiley, 2004, Chapters~4 and~6. Sylow theory with many worked classifications.
\item R.~Gow, \emph{Sylow's proof of Sylow's theorem}, Irish Math.\ Soc.\ Bulletin 33 (1994), 55--63. An account quoting the original statements and comparing their proof with later approaches.
\item \emph{Historical biography:} MacTutor History of Mathematics, \emph{Ludwig Sylow}, \url{https://mathshistory.st-andrews.ac.uk/Biographies/Sylow/}. The original 1872 paper is cited above.
\end{itemize}
